\section{Propiedades matemáticas y verificación formal}\label{sec:propiedades}

El programa \eqref{eq:objetivo_smcfhr} admite una caracterización analítica completa en regularidad diferencial, convexidad, existencia, unicidad, condiciones de optimalidad y sensibilidad (demostraciones completas, matriz Hessiana por tramos y análisis de complejidad en el Material Suplementario, Sección~S1 y Figura~S6).

\subsection{Gradiente analítico y regularidad $C^1$}

La derivada de Huber $\psi_\delta(r) := \rho_\delta'(r) = \clip(r, -\delta, \delta)$ es continua y $1$-Lipschitz en $\mathbb{R}$. Definiendo el vector de residuos $\mathbf{r}(b, \mathbf{w}) := \tilde{\mathbf{y}} - b\mathbf{1}_n - Z\mathbf{w} \in \mathbb{R}^n$, el gradiente exacto de $J$ es:
\begin{equation}
\frac{\partial J}{\partial b}(b,\mathbf{w}) = -\frac{1}{n}\sum_{k=1}^{n} \psi_\delta(r_k), \qquad
\nabla_{\mathbf{w}} J(b,\mathbf{w}) = -\frac{1}{n} Z^\top \psi_\delta(\mathbf{r}) + D_\alpha \mathbf{w}.
\label{eq:gradiente_smcfhr}
\end{equation}
La función objetivo $J$ es de clase $C^1(\mathbb{R}^{p+1})$ con gradiente Lipschitz continuo, pero no es $C^2$ sobre los hiperplanos $|r_k(b,\mathbf{w})| = \delta$, donde la segunda derivada presenta saltos finitos. Esta regularidad $C^1$ por tramos justifica el uso de L-BFGS-B \cite{byrd1995lbfgsb,zhu1997lbfgsb}: emplea únicamente evaluaciones de $J$, gradientes analíticos de primer orden y proyecciones simples sobre $\mathbb{R} \times \mathcal{C}$ con coste $O(np)$ por evaluación, evitando invertir Hessianas discontinuas.

\subsection{Convexidad, existencia, unicidad, KKT e influencia}

Distinguiendo las consecuencias directas de la teoría convexa \cite{boyd2004convex} y estimación $M$ \cite{huber1964,huber2009robust} frente a las propiedades específicas de un intercepto no penalizado bajo pérdida de Huber:

\begin{theorem}[Convexidad y existencia de soluciones]\label{thm:convexidad_existencia}
Si $n \ge 1$, $\delta > 0$, $\alpha_{\mathrm{lin}} > 0$ y $\alpha_{\mathrm{meteo}} > 0$, entonces: (1) $\Omega := \mathbb{R} \times \mathcal{C}$ es cerrado y convexo en $\mathbb{R}^{p+1}$; (2) $J \in C^1(\mathbb{R}^{p+1})$ es convexa con gradiente globalmente Lipschitz de constante $L_J \le \frac{1}{n}\left\|\begin{bmatrix}\mathbf{1}_n & Z\end{bmatrix}\right\|_2^2 + \max\{\alpha_{\mathrm{lin}}, \alpha_{\mathrm{meteo}}\}$; y (3) $J$ es coerciva sobre $\Omega$ y acotada inferiormente por $0$, garantizando al menos un minimizador global $(b^*, \mathbf{w}^*) \in \Omega$.
\end{theorem}

\begin{theorem}[Unicidad exacta y caso degenerado del intercepto]\label{thm:unicidad}
Bajo las hipótesis del Teorema~\ref{thm:convexidad_existencia}: (1) el vector óptimo de coeficientes $\mathbf{w}^* \in \mathcal{C}$ es \textbf{único}; (2) si existe al menos una observación $k$ con $|\tilde{y}_k - b^* - \mathbf{z}_k^\top \mathbf{w}^*| < \delta$ para algún $b^*$ óptimo, entonces el intercepto $b^*$ también es \textbf{único}; y (3) si $|r_k(b^*,\mathbf{w}^*)| \ge \delta$ para todo $k$ (con $n$ par e igual número de residuos positivos y negativos), el conjunto de interceptos óptimos puede ser un intervalo no degenerado.
\end{theorem}

Dado que $\mathcal{C}$ está definido por desigualdades afines y satisface la condición de Slater ($w_j = 1 > 0$ para $j \in \mathcal{M}_+$), las condiciones KKT son necesarias y suficientes para la optimalidad global (Material Suplementario, Sección~S1, Proposición~S1): en particular, cuando una cota $w_j^* = 0$ ($j \in \mathcal{M}_+$) está activa, la correlación residual ponderada cumple $\frac{1}{n}(Z^\top \psi_\delta(\mathbf{r}^*))_j \le 0$. Asimismo, aunque $|\psi_\delta(r_k)| \le \delta$ acota la influencia de cada observación en la respuesta ($\delta/n$ en el intercepto y $(\delta/n)\|\mathbf{z}_k\|_\infty$ en los pesos), SM-CFHR \emph{no} posee función de influencia acotada frente a puntos de apalancamiento arbitrarios en covariables $\|\mathbf{z}_k\|$ \cite{hampel1986robust}. Por último, para todo $j \in \mathcal{M}_+$ con $\tilde{s}_j > 0$, $\hat{y}_q$ es monótona no decreciente en $x_{q,j}$ \textit{ceteris paribus} y permanece en $[0, 24\,000\,H_{\mathrm{dir}}\,C_{i,t}]$, sin garantizar monotonía conjunta si múltiples covariables correlacionadas varían simultáneamente.

\subsection{Verificación formal en Lean~4}\label{sec:lean4_resumen}

Los resultados analíticos anteriores fueron formalizados en \textbf{Lean~4} (\texttt{leanprover/lean4:v4.36.0}) \cite{moura2021lean4} sobre \textbf{Mathlib} \cite{mathlib2020} (Material Suplementario, Sección~S2, Tabla~S1), compilando $26$ declaraciones verificadas en seis módulos sobre $\mathbb{R}$ sin axiomas personalizados (únicamente \texttt{propext}, \texttt{Classical.choice} y \texttt{Quot.sound}). La formalización distingue hipótesis explícitas ($n \ge 1$ y $\alpha_j > 0$ para coercividad/existencia frente a $\alpha_j \ge 0$ para convexidad/KKT) y demuestra la regularidad Lipschitz de $\partial J/\partial b$ (\texttt{gradB\_lipschitz}), mientras la cota espectral completa $L_J$ de $\nabla J$ se demuestra analíticamente en la Sección~S1. La verificación en $\mathbb{R}$ no cubre aritmética IEEE-754 ni convergencia numérica de L-BFGS-B.
